\chapter{Introduction}

The aim of this work is to present a study concerning the development of an
expert system, built in a rule-based programming environment, for the
recommendation of a beer style according to the user's taste and to the meal to
be paired.

Throughout the report we document the various design choices adopted during the
definition phase, as well as the behaviour that the expert system exhibits
during its execution.

Before discussing the project in detail, we provide a brief introduction to the
features that characterise an expert system, to expert systems applied to the
recommendation task, and to the handling of uncertain knowledge.


\section{Expert Systems}

One of the most actively developed areas in Artificial Intelligence is the study
and design of expert systems, also referred to as \emph{Expert Systems} or
\emph{Knowledge-Based Systems}. In the following sections we illustrate the
features and the components that characterise an expert system.


\subsection{Definition}

The term \emph{expert system} was introduced in 1977 by Feigenbaum, who also
provided a precise definition:

\begin{quotation}
An expert system is a computer program that uses knowledge and reasoning
techniques to solve problems that would normally require the help of a human
expert. An expert system must be able to justify or explain why a particular
solution was reached for a given problem.
\end{quotation}

An expert system must be capable of emulating the work of a human expert in a
given domain. In particular, an expert system must be able to perform the same
actions, provide the same judgements, and offer the same explanations as a
human. This feature, however, should not obscure the ultimate goal of an expert
system. Indeed, contrary to what one might think, its purpose is not to replace
the expert, but rather to assist and help the user improve their performance.

A distinctive trait of this paradigm is the clear separation between
\emph{knowledge} ---the facts and rules that describe the domain--- and the
\emph{inference} mechanism that applies them. Conventional programming is
designed for procedural data manipulation, whereas humans solve many problems
through abstract, symbolic reasoning. Expert systems capture such reasoning by
encoding it declaratively, so that the body of knowledge can be inspected,
maintained and extended independently of the control flow.

The fundamental features of an expert system are the following:

\begin{itemize}
    \item \emph{Possession of specific knowledge}: it is important that an expert
    system possesses enough domain-specific knowledge. The quality of an expert
    system, in fact, does not depend only on the formalism or the inference
    technique adopted, but also on the knowledge that the program possesses;

    \item \emph{Adherence of the domain to reality}: the domain in which an
    expert system operates must not be a simplified model of the real domain, but
    must represent it as it is, of course within certain limits, yet never
    resorting to excessively abstract models;

    \item \emph{Reasoning capability}: an expert system must draw inferences and
    make decisions by exploiting the available knowledge base;

    \item \emph{Capability of interacting with the external world}: an expert
    system must have a communication protocol with the external world. Such
    communication is necessary in order to retrieve the information required to
    carry on the reasoning and to provide explanations to the user;

    \item \emph{Capability of explaining its own operation}: a good expert system
    must be able to make explicit and explain the decisions taken and the
    reasoning by which it reached a conclusion. This capability strengthens the
    degree of trust that the user places in the system;

    \item \emph{Learning capability}: a feature that, while not common to all
    expert systems, is very interesting to consider. Just as a human expert
    learns new experiences and broadens their knowledge over time, an expert
    system should likewise be able to learn new knowledge.
\end{itemize}

A widely adopted way of representing declarative knowledge is the
\emph{production rule}, an \code{if-then} construct whose left-hand side states
the conditions under which the right-hand side actions must be taken. Production
systems are typically driven by \emph{forward chaining}: the engine starts from
the known facts and repeatedly fires the rules whose conditions are satisfied,
deriving new facts until no further rule applies. CLIPS, the environment adopted
in this work, implements this strategy efficiently by means of the
\emph{Rete} pattern-matching algorithm, which incrementally tracks which rules
match the current contents of the working memory and thus avoids re-evaluating
every rule on each cycle.


\subsection{Components}

Within an expert system it is possible to identify four fundamental components:
the knowledge base, the inference engine, the user interface, and the working
memory.

\subsubsection*{Knowledge Base}

The knowledge base represents the set of information needed to address and solve
a problem in a specific domain.

Such information is called \emph{knowledge}, and it must be represented by means
of a formalism that the machine can interpret (for example, Horn logic,
description logic, modal logic, and so on).

Typically, the knowledge base is organised into two blocks:

\begin{itemize}
    \item \emph{Block of assertions or facts}: it represents the declarative
    knowledge related to a particular problem to be solved, that is, true facts
    introduced at the beginning of the consultation. In CLIPS the shape of a
    fact is declared through a \code{deftemplate}, while sets of initial facts
    are conveniently grouped within a \code{deffacts} construct;

    \item \emph{Block of relations or rules}: it contains the relations among the
    elementary facts that are useful for building the body of knowledge. These
    are usually expressed in the form of production rules ---the \code{defrule}
    construct in CLIPS--- or according to other types of representation such as
    semantic networks, frames, scripts, and so on.
\end{itemize}

\subsubsection*{Inference Engine}

The inference engine is a program able to deduce, starting from the knowledge
base, the conclusions that constitute the solution of the problem in the domain
in which the expert system operates.

The inference engine is able to identify the set of rules that are applicable at
a given moment, to select the rule to apply among the candidate ones, and to
execute the action indicated by the selected rule. In CLIPS the rules whose
conditions are satisfied are placed on the \emph{agenda}; when more than one rule
is eligible, a \emph{conflict resolution} strategy (based, for instance, on rule
salience) decides the order in which they are fired.

\subsubsection*{User Interface}

The user interface is responsible for facilitating the interaction between the
user and the expert system. It can be built using user-friendly graphical
interfaces or command-line interfaces. It is important that the interface makes
the dialogue with the user as close as possible to a natural one.

\subsubsection*{Working Memory}

The working memory is the component that allows the storage of newly observed
evidence or of data useful for solving the problem during a specific working
session. The content of this memory will therefore change throughout the entire
execution session of the expert system, as facts are asserted and retracted in
response to the user's answers and to the firing of rules.


\subsection{Roles}

In general, the roles that interact with an expert system are three:

\begin{itemize}
    \item \emph{Domain expert}: the actor who holds the knowledge and experience
    about a specific domain and is able to solve the problem at hand ---in our
    case, the sommelier-like expertise that relates beer styles to tastes,
    occasions and food;

    \item \emph{Knowledge engineer}: the actor who encodes the knowledge,
    provided by the domain expert, into a formalism that the expert system can
    understand;

    \item \emph{User}: the individual who consults the expert system in order to
    benefit from it.
\end{itemize}


\subsection{Building Environments}

The construction of an expert system can be carried out using two different types
of environment:

\begin{itemize}
    \item \emph{Skeleton environments}: these are environments derived from
    existing expert systems, from which the knowledge related to the application
    domain has been removed. Examples include Emycin, Kas, Expert, and so on;

    \item \emph{Environment-type environments}: these are environments that
    provide a language for representing knowledge. Examples include Clips,
    Prolog, and so on.
\end{itemize}

Comparing the two types of environment, one can state that ``Skeleton'' systems
are more immediate, but less flexible, since they are affected by the
specialised architecture of the original product; ``Environment''-type systems,
on the other hand, have flexibility as their strength, while their weakness is
the greater complexity.

BeerEX is built on CLIPS, a forward-chaining production-system tool that acts as
a general-purpose shell for rule-based reasoning. Because knowledge is kept
declarative, new styles, questions or food pairings can be added without
touching the control flow, which is a decisive advantage for the maintainability
of the system.


\section{Expert Systems for Recommendation}

A second field, beyond the classical diagnostic ones, in which the expert-system
paradigm proves effective is that of \emph{recommendation} and advisory tasks,
and in particular within the food-and-beverage domain.

The craft-beer market is vast and fast-growing: in the United States alone there
are thousands of breweries and well over a hundred recognised beer styles, with
tens of thousands of brands. For the consumer this abundance translates into
uncertainty: choosing a beer that suits one's taste, the season, or ---above
all--- the food on the table requires expertise that most drinkers do not have.
BeerEX addresses exactly this gap, acting as a pocket beer sommelier that, by
asking a few simple questions, recommends in real time the most appropriate beer
\emph{style} for the situation, together with an explanation of the choice.

Such a task is a natural fit for an expert system. The domain is a bounded,
symbolic one: the styles, their qualitative and serving profiles and the
recommended food pairings can be enumerated and described declaratively. The
knowledge involved is genuine expertise, comparable to that of a sommelier,
which a non-expert user typically lacks. Finally, a recommendation is far more
valuable when it is accompanied by an \emph{explanation}: an expert system can
make explicit \emph{why} a question is asked and \emph{how}, factor by factor, a
particular style came to be suggested, thereby strengthening the user's trust in
the advice. As with other advisory expert systems, the consultation allows a
non-expert person to reach a decision that would otherwise lie beyond their level
of knowledge and experience, while preserving and reusing the encoded expertise
consistently across every consultation.


\section{Handling Uncertain Knowledge}

The cognitive mechanism of the human expert often relies on empirical rules
acquired from their own experience and on data that are not completely certain.

An expert is able to reason even when the available information is uncertain or
incomplete. This implies that conclusions obtained from knowledge held to be
generally true, and by means of approximate reasoning, are not always absolutely
correct.

This is precisely the case for a beer recommendation, which is rarely a
black-or-white matter: tastes are graded, occasions overlap, and a given style
may suit a dish more or less well rather than perfectly or not at all. The need
therefore emerges for approaches that are able to replicate this approximate
reasoning within the expert system.

There are several ways to represent and formalise uncertainty: they range from
quantitative methods (certainty factors, Bayesian networks, belief functions,
the fuzzy approach) to symbolic methods (non-monotonic logics). All these
methods share the common feature of having to deal with the combination of
evidence and the propagation of uncertainty from facts to conclusions.

BeerEX offers two interchangeable calculi for reasoning under uncertainty. The
first is the classic \emph{certainty-factor} model introduced by MYCIN, in which
each piece of evidence and each conclusion carries a degree of belief, a \cf,
and the rules combine and propagate these factors. The second is a
\emph{fuzzy-logic} engine whose membership functions are data-grounded, that is,
derived from measurable beer parameters ---alcohol content, colour and
bitterness--- so that linguistic judgements correspond to real style data. The
user (or the deployer) can switch between the two backends, and the two can be
compared on the same session. The details of both calculi are deferred to later
chapters.


\section{Scope of this report}

The report follows the knowledge-engineering lifecycle: the analysis of the
domain (Chapter~2), the identification of the task (Chapter~3), the
conceptualization of the knowledge and of the two uncertainty models
(Chapter~4), the CLIPS implementation and the Telegram delivery (Chapter~5), the
ways to run the system (Chapter~6) and its validation against an authentic
FuzzyCLIPS engine (Chapter~7).
